\documentclass[aspectratio=169,utf8,t]{beamer}

% Idioma e Codificacao
\usepackage[brazilian]{babel}
\usepackage{booktabs}
\usepackage{array}
\usepackage[normalem]{ulem}
\usepackage{tikz}
\usetikzlibrary{shapes,shapes.geometric,arrows,arrows.meta,positioning,calc,fit}

% Carrega o Tema Modelo dos Scouts (Widescreen 16:9, Nunito Sans)
\usetheme{scouts}

\definecolor{amber}{RGB}{245, 158, 11}

% Metadados Institucionais
\author{Prof. Me. Reginaldo Fernandes}
\institute{Instituto Federal do Ceará -- Campus Tauá \\ Tecnologia em Telemática}
\title{Banco de Dados}
\subtitle{Aula 09: Fundamentos Formais do Modelo Relacional e Restrições de Integridade}
\date{21 de Outubro de 2026}

\begin{document}

% ------------------------------------------------------------------------------
% Capa da Apresentacao
% ------------------------------------------------------------------------------
\changecolours[inverse]{ScoutGreen}
\frame[plain]{\titlepage}

% ------------------------------------------------------------------------------
% Sumario Executivo da Aula
% ------------------------------------------------------------------------------
\changecolours{ScoutGreen}
\begin{frame}{Agenda do Encontro}
  \begin{columns}[t]
    \begin{column}{0.48\textwidth}
      \scriptsize
      \tableofcontents[sections={1-5}]
    \end{column}
    \begin{column}{0.48\textwidth}
      \scriptsize
      \tableofcontents[sections={6-9}]
    \end{column}
  \end{columns}
\end{frame}

% ------------------------------------------------------------------------------
% Modulos Tematicos da Aula
% ------------------------------------------------------------------------------
\changecolours{ScoutGreen}
\section[Revisão e Objetivos]{A Revolução de Codd (1970) e Objetivos da Aula}

% ------------------------------------------------------------------------------
\twocolframe[title={Do DER Conceitual (Aula 08) ao Modelo Relacional (Aula 09)},leftcol=7.0cm,rightcol=7.0cm]{%
  \alert{A Revolução de Edgar Frank Codd (1970)}\par
  \itemseps[topsepi=0.5mm,itemsepi=1mm,labelsep=2mm,leftmargini=3mm]
  \begin{itemize}
    \item \textbf{Artigo Histórico (IBM):} \textit{``A Relational Model of Data for Large Shared Data Banks''} (CACM, Junho de 1970).
    \item \textbf{Ruptura com Ponteiros Físicos:} Antes de 1970, bancos hierárquicos e em rede exigiam navegar por endereços físicos de disco.
    \item \textbf{Base Matemática Sólida:} Todos os dados são representados exclusivamente como \textbf{relações matemáticas} baseadas na teoria dos conjuntos.
  \end{itemize}
}{%
  \alert{A Transição no Projeto de BD}\par
  \itemseps[topsepi=0.5mm,itemsepi=1mm,labelsep=2mm,leftmargini=3mm]
  \begin{itemize}
    \item \textbf{Nível Conceitual (Aulas 03 a 08):} Focado na semântica do negócio (Entidades, Losangos, Herança EER).
    \item \textbf{Nível Lógico Relacional (Aula 09 em diante):} Focado na representação matemática e tabular de \textbf{Relações, Tuplas, Chaves e Restrições}.
    \item \textbf{Meta da Aula 09:} Dominar os fundamentos formais antes de mapear o DER em esquemas relacionais (Aula 10).
  \end{itemize}
}

% ------------------------------------------------------------------------------
\twocolframe[title={Objetivos de Aprendizagem da Aula 09},leftcol=7.0cm,rightcol=7.0cm]{%
  \alert{Competências Teóricas e Matemáticas}\par
  \itemseps[topsepi=0.5mm,itemsepi=1mm,labelsep=2mm,leftmargini=3mm]
  \begin{itemize}
    \item \textbf{Anatomia da Relação:} Definir Domínio ($D_i$), Atributo ($A_i$), Esquema $R$, Tupla $t$ e Instância $r(R)$.
    \item \textbf{Métricas Estruturais:} Diferenciar \textbf{Grau (Aridade)} de \textbf{Cardinalidade}.
    \item \textbf{Propriedades de Conjuntos:} Atomicidade (1FN), ausência de ordem de tuplas e proibição de duplicatas.
    \item \textbf{Teoria das Chaves:} Superchave ($SK$), Chave Candidata ($CK$), Primária ($PK$) e Alternativa ($AK$).
  \end{itemize}
}{%
  \alert{Competências de Integridade e Modelagem}\par
  \itemseps[topsepi=0.5mm,itemsepi=1mm,labelsep=2mm,leftmargini=3mm]
  \begin{itemize}
    \item \textbf{As 4 Restrições Formais:} Domínio, Chave, Entidade ($\text{PK} \neq \text{NULL}$) e Referencial ($FK$).
    \item \textbf{Semântica do Valor Nulo:} Compreender quando a chave estrangeira aceita valor nulo e quando a PK o proíbe.
    \item \textbf{Políticas Referenciais:} Avaliar Restrição, Propagação em Cascata e Anulação na exclusão de tuplas.
    \item \textbf{Diagrama de Esquema:} Desenhar relações e arcos direcionados de integridade referencial.
  \end{itemize}
}

\changecolours{ScoutGreen}
\section[Anatomia da Relação]{Domínios, Atributos, Tuplas, Grau e Cardinalidade}

% ------------------------------------------------------------------------------
\begin{frame}{Dicionário de Equivalência Terminológica e Anatomia da Relação}
  \vspace{-0.1cm}
  \begin{columns}[t]
    \begin{column}{0.46\textwidth}
      \alert{Três Vocabulários Equivalentes}\par
      \vspace{0.1cm}
      \scriptsize
      \begin{tabular}{lll}
        \toprule
        \textbf{Formal (Codd)} & \textbf{Tabular / Lógico} & \textbf{Arquivos} \\
        \midrule
        Relação ($R$) & Tabela & Arquivo \\
        Tupla ($t$) & Linha & Registro \\
        Atributo ($A_i$) & Coluna & Campo \\
        Domínio ($D_i$) & Tipo / Formato & Máscara \\
        Grau ($n$) & Nº de Colunas & Nº Campos \\
        Cardinalidade ($m$) & Nº de Linhas & Nº Registros \\
        \bottomrule
      \end{tabular}
      \vspace{0.15cm}\par
      $\bullet$ \textbf{Grau (Aridade):} Propriedade \textbf{estática} do esquema $R$.\\
      $\bullet$ \textbf{Cardinalidade:} Propriedade \textbf{dinâmica} do estado $r(R)$ (varia conforme novas tuplas entram ou saem).
    \end{column}
    \begin{column}{0.51\textwidth}
      \alert{Visualização Formal: Relação \texttt{MEDICO}}\par
      \vspace{0.1cm}
      \begin{center}
        \tikz[scale=0.68, transform shape]{
          \node[font=\bfseries\sffamily\scriptsize, ScoutNavy] at (2.7, 2.1) {Esquema $R$: \texttt{MEDICO}(\underline{matricula}, nome, crm, cod\_depto)};

          % Cabecalho (Atributos)
          \draw[fill=ScoutGreen!20, draw=ScoutGreen, thick] (0, 1.1) rectangle (5.4, 1.7);
          \node[font=\bfseries\sffamily\tiny] at (0.65, 1.4) {\underline{matricula}};
          \node[font=\bfseries\sffamily\tiny] at (2.05, 1.4) {nome};
          \node[font=\bfseries\sffamily\tiny] at (3.45, 1.4) {crm};
          \node[font=\bfseries\sffamily\tiny] at (4.75, 1.4) {cod\_depto};

          % Tuplas
          \draw[fill=white, draw=black!50] (0, 0.5) rectangle (5.4, 1.1);
          \node[font=\sffamily\tiny] at (0.65, 0.8) {1001};
          \node[font=\sffamily\tiny] at (2.05, 0.8) {Dra. Ana};
          \node[font=\sffamily\tiny] at (3.45, 0.8) {CE-1234};
          \node[font=\sffamily\tiny] at (4.75, 0.8) {10};

          \draw[fill=black!3, draw=black!50] (0, -0.1) rectangle (5.4, 0.5);
          \node[font=\sffamily\tiny] at (0.65, 0.2) {1002};
          \node[font=\sffamily\tiny] at (2.05, 0.2) {Dr. Bruno};
          \node[font=\sffamily\tiny] at (3.45, 0.2) {CE-5678};
          \node[font=\sffamily\tiny] at (4.75, 0.2) {10};

          \draw[fill=white, draw=black!50] (0, -0.7) rectangle (5.4, -0.1);
          \node[font=\sffamily\tiny] at (0.65, -0.4) {1003};
          \node[font=\sffamily\tiny] at (2.05, -0.4) {Dra. Carla};
          \node[font=\sffamily\tiny] at (3.45, -0.4) {CE-9012};
          \node[font=\sffamily\tiny\itshape, gray] at (4.75, -0.4) {NULO};

          % Cotas Grau e Cardinalidade
          \draw[<->, thick, ScoutNavy] (0, -1.1) -- node[below, font=\tiny\bfseries] {Grau (Aridade) $n = 4$ atributos} (5.4, -1.1);
          \draw[<->, thick, ScoutGreen] (5.7, -0.7) -- node[right, font=\tiny\bfseries, align=left] {Cardinalidade\\$m = 3$ tuplas} (5.7, 1.1);
        }
      \end{center}
    \end{column}
  \end{columns}
\end{frame}

% ------------------------------------------------------------------------------
\twocolframe[title={Definição Matemática de Esquema ($R$) e Instância ($r(R)$)},leftcol=7.0cm,rightcol=7.0cm]{%
  \alert{Domínio Atômico e Esquema Intensional}\par
  \itemseps[topsepi=0.5mm,itemsepi=1mm,labelsep=2mm,leftmargini=3mm]
  \begin{itemize}
    \item \textbf{Domínio ($D_i$):} Conjunto de valores \textbf{atômicos} (indivisíveis) permitidos para um atributo, associado a uma definição semântica e formato.
    \item \textbf{Esquema de Relação ($R$):} Denotado por $R(A_1, A_2, \dots, A_n)$, define a estrutura intensional (cabeçalho invariante) composta por $n$ atributos onde $\text{dom}(A_i) = D_i$.
  \end{itemize}
}{%
  \alert{Estado/Instância como Subconjunto Cartesiano}\par
  \itemseps[topsepi=0.5mm,itemsepi=1mm,labelsep=2mm,leftmargini=3mm]
  \begin{itemize}
    \item \textbf{Instância Extensional $r(R)$:} É um conjunto finito de $m$ tuplas $\{t_1, t_2, \dots, t_m\}$ que reflete o estado dos dados em dado instante:
      $$r(R) \subseteq \text{dom}(A_1) \times \text{dom}(A_2) \times \dots \times \text{dom}(A_n)$$
    \item \textbf{Cada Tupla $t$:} É uma lista ordenada de $n$ valores $t = \langle v_1, v_2, \dots, v_n \rangle$, onde cada $v_i \in \text{dom}(A_i)$ ou é um valor nulo especial (\textit{NULL}).
  \end{itemize}
}

\changecolours{ScoutGreen}
\section[Propriedades e NULL]{Propriedades Matemáticas das Relações e o Valor Nulo}

% ------------------------------------------------------------------------------
\twocolframe[title={As 4 Propriedades Fundamentais das Relações e o Valor Nulo},leftcol=7.0cm,rightcol=7.0cm]{%
  \alert{Consequências de $r(R)$ Ser um Conjunto}\par
  \itemseps[topsepi=0.5mm,itemsepi=1mm,labelsep=2mm,leftmargini=3mm]
  \begin{itemize}
    \item \textbf{1. Atomicidade de Valores (1FN):} São proibidos atributos compostos ou multivalorados na mesma célula. Cada posição da tupla guarda um único valor indivisível.
    \item \textbf{2. Inexistência de Ordem das Tuplas:} Um conjunto matemático não possui ordem ($\{t_1, t_2\} = \{t_2, t_1\}$).
    \item \textbf{3. Proibição de Tuplas Duplicadas:} Em um conjunto não existem dois elementos idênticos. Toda relação deve possuir chave primária!
  \end{itemize}
}{%
  \alert{A Semântica Formal do Valor Nulo (\textit{NULL})}\par
  \itemseps[topsepi=0.5mm,itemsepi=1mm,labelsep=2mm,leftmargini=3mm]
  \begin{itemize}
    \item \textbf{Conceito:} Não representa o número zero nem texto em branco; indica ausência ou desconhecimento de valor aplicável.
    \item \textbf{1. Valor Desconhecido:} Existe no mundo real, mas não é sabido (ex: data de nascimento não informada).
    \item \textbf{2. Valor Inaplicável:} O atributo não se aplica àquela tupla específica.
    \item \textbf{3. Valor Indisponível:} Existe, mas ainda não foi coletado no momento do registro.
  \end{itemize}
}

\changecolours{ScoutGreen}
\section[Teoria das Chaves]{Superchave, Chave Candidata, Primária e Alternativa}

% ------------------------------------------------------------------------------
\twocolframe[title={Hierarquia Matemática de Chaves: $SK \supset CK \supset PK$},leftcol=7.0cm,rightcol=7.0cm]{%
  \alert{Superchave ($SK$) vs. Chave Candidata ($CK$)}\par
  \itemseps[topsepi=0.5mm,itemsepi=1mm,labelsep=2mm,leftmargini=3mm]
  \begin{itemize}
    \item \textbf{Superchave (\textit{Superkey} -- $SK$):} Qualquer subconjunto de atributos $SK \subseteq R$ que garante unicidade: para quaisquer duas tuplas distintas $t_1 \neq t_2$, vale $t_1[SK] \neq t_2[SK]$.
    \item \textbf{Exemplos em \texttt{MEDICO}:} $\{\text{matricula}\}$, $\{\text{matricula}, \text{nome}\}$ e $\{\text{cpf}, \text{nome}\}$ são todas superchaves!
    \item \textbf{Chave Candidata ($CK$ -- Minimalidade):} É uma \textbf{superchave mínima (irredutível)} --- se removermos qualquer atributo dela, ela deixa de ser única!
  \end{itemize}
}{%
  \alert{Chave Primária ($PK$) e Alternativa ($AK$)}\par
  \itemseps[topsepi=0.5mm,itemsepi=1mm,labelsep=2mm,leftmargini=3mm]
  \begin{itemize}
    \item \textbf{Minimalidade em \texttt{MEDICO}:} $\{\text{matricula}\}$, $\{\text{cpf}\}$ e $\{\text{crm}\}$ são Chaves Candidatas ($CK$). Já $\{\text{matricula}, \text{nome}\}$ é apenas Superchave (pois \texttt{nome} é supérfluo para unicidade).
    \item \textbf{Chave Primária ($PK$):} A chave candidata eleita pelo projetista como identificador principal ($\underline{\text{matricula}}$).
    \item \textbf{Chaves Alternativas ($AK$):} As demais chaves candidatas não escolhidas como PK (\texttt{cpf}, \texttt{crm}).
  \end{itemize}
}

\changecolours{ScoutGreen}
\section[Domínio e Entidade]{Restrições de Domínio, Chave e Integridade de Entidade}

% ------------------------------------------------------------------------------
\twocolframe[title={As 3 Primeiras Restrições Formais do Modelo Relacional},leftcol=7.0cm,rightcol=7.0cm]{%
  \alert{1. Domínio e 2. Unicidade de Chave}\par
  \itemseps[topsepi=0.5mm,itemsepi=1mm,labelsep=2mm,leftmargini=3mm]
  \begin{itemize}
    \item \textbf{1. Restrição de Domínio (e 1FN):} Em toda tupla $t$, o valor de cada atributo $t[A_i]$ deve pertencer ao domínio atômico $\text{dom}(A_i)$, proibindo valores multivalorados ou tipos incompatíveis.
    \item \textbf{2. Restrição de Chave (Unicidade):} Duas tuplas distintas não podem ter valores idênticos na Chave Primária ($PK$) nem nas Chaves Candidatas Alternativas ($AK$).
  \end{itemize}
}{%
  \alert{3. Regra da Integridade de Entidade (Codd)}\par
  \itemseps[topsepi=0.5mm,itemsepi=1mm,labelsep=2mm,leftmargini=3mm]
  \begin{itemize}
    \item \textbf{Enunciado Formal:} \textit{``Nenhum atributo que participe da Chave Primária ($PK$) de uma relação base pode assumir valor nulo (\texttt{NULL}).''}
      $$\forall t \in r(R), \quad t[PK] \neq \text{NULL}$$
    \item \textbf{Justificativa Lógica:} A PK identifica cada tupla individualmente. Se $t[PK]$ fosse nulo, seria impossível distinguir tuplas ou referenciá-las a partir de outras relações!
  \end{itemize}
}

\changecolours{ScoutGreen}
\section[Integridade Referencial]{Chave Estrangeira (FK) e Integridade Referencial}

% ------------------------------------------------------------------------------
\begin{frame}{4. Restrição de Integridade Referencial e Chave Estrangeira (FK)}
  \vspace{-0.1cm}
  \begin{center}
    \tikz[scale=0.72, transform shape]{
      % Relacao Referenciada (Destino)
      \node[font=\bfseries\sffamily\scriptsize, ScoutGreen] at (-4.2, 1.15) {Relação Referenciada ($R_2$ -- Pai):};
      \node[draw=ScoutGreen, fill=ScoutGreen!15, thick, minimum width=5.2cm, minimum height=0.65cm, font=\sffamily\scriptsize] (r2) at (-4.2, 0.35) {
        \textbf{DEPARTAMENTO\_CLINICO}(\underline{cod\_depto}, nome\_depto, bloco)
      };

      % Relacao Referenciadora (Origem)
      \node[font=\bfseries\sffamily\scriptsize, ScoutNavy] at (3.6, 1.15) {Relação Referenciadora ($R_1$ -- Filha):};
      \node[draw=ScoutNavy, fill=ScoutNavy!12, thick, minimum width=6.6cm, minimum height=0.65cm, font=\sffamily\scriptsize] (r1) at (3.6, 0.35) {
        \textbf{MEDICO}(\underline{matricula}, cpf, nome, crm, \textbf{\color{ScoutNavy}cod\_depto\_fk})
      };

      % Arco Direcionado da FK para a PK
      \draw[-{Latex[length=2.8mm]}, very thick, ScoutNavy] (5.9, 0.0) -- (5.9, -0.65) -- (-4.7, -0.65) -- (-4.7, 0.0);
      \node[font=\tiny\bfseries\sffamily, ScoutNavy] at (0.6, -0.95) {Arco de Integridade Referencial: $\text{MEDICO}[\text{cod\_depto\_fk}] \longrightarrow \text{DEPARTAMENTO\_CLINICO}[\underline{\text{cod\_depto}}]$};
    }
  \end{center}
  \vspace{0.05cm}
  \begin{columns}[t]
    \begin{column}{0.49\textwidth}
      \footnotesize
      \alert{Regra Matemática da FK (Codd):}\par
      Uma tupla $t_1 \in r(R_1)$ satisfaz a integridade referencial se, e somente se:\par
      \vspace{0.1cm}
      \textbf{(a)} $\exists t_2 \in r(R_2)$ tal que $t_1[FK] = t_2[PK]$, \textbf{OU}\\
      \textbf{(b)} $t_1[FK] = \text{NULL}$ (quando a participação for parcial/opcional e a FK não compuser a PK de $R_1$).
    \end{column}
    \begin{column}{0.48\textwidth}
      \footnotesize
      \alert{Políticas de Manutenção na Exclusão de $R_2$:}\par
      $\bullet$ \textbf{Restrição (\textit{Restrict}):} Proíbe excluir o departamento se houver médicos vinculados a ele.\\
      $\bullet$ \textbf{Anulação (\textit{Set Null}):} Exclui o departamento e define a FK dos médicos como \texttt{NULL}.\\
      $\bullet$ \textbf{Cascata (\textit{Cascade}):} Propaga a exclusão para as tuplas dependentes (usado em entidades fracas).
    \end{column}
  \end{columns}
\end{frame}

\changecolours{ScoutGreen}
\section[Auditoria Tabular]{Estudo de Caso e Auditoria Tabular de Instâncias Relacionais}

% ------------------------------------------------------------------------------
\begin{frame}{Diagrama de Esquema Completo: Clínica Médica ``VidaPlena''}
  \vspace{-0.22cm}
  \begin{center}
    \tikz[scale=0.71, transform shape]{
      % 1. DEPARTAMENTO_CLINICO
      \node[draw=ScoutGreen, fill=ScoutGreen!15, thick, minimum width=6.1cm, minimum height=0.65cm, font=\sffamily\scriptsize] (depto) at (-4.3, 1.45) {
        \textbf{DEPARTAMENTO\_CLINICO}(\underline{cod\_depto}, nome\_depto, bloco)
      };

      % 2. MEDICO
      \node[draw=ScoutNavy, fill=ScoutNavy!12, thick, minimum width=8.1cm, minimum height=0.65cm, font=\sffamily\scriptsize] (med) at (4.1, 1.45) {
        \textbf{MEDICO}(\underline{matricula}, cpf, nome, crm, \textbf{\color{ScoutNavy}cod\_depto\_fk}, \textbf{\color{ScoutGreen}mat\_superv\_fk})
      };

      % 3. PACIENTE
      \node[draw=ScoutGreen, fill=ScoutGreen!15, thick, minimum width=6.1cm, minimum height=0.65cm, font=\sffamily\scriptsize] (pac) at (-4.3, -0.50) {
        \textbf{PACIENTE}(\underline{id\_prontuario}, cpf, nome\_paciente, data\_nasc)
      };

      % 4. DEPENDENTE (PK composta com FK)
      \node[draw=ScoutNavy, fill=amber!15, thick, minimum width=7.6cm, minimum height=0.65cm, font=\sffamily\scriptsize] (dep) at (4.1, -0.50) {
        \textbf{DEPENDENTE}(\underline{\textbf{\color{ScoutNavy}id\_prontuario\_fk}, \dashuline{nome\_dep}}, data\_nasc, parentesco)
      };

      % Arco 1: MEDICO.cod_depto_fk -> DEPARTAMENTO_CLINICO.cod_depto
      \draw[-{Latex[length=2.2mm]}, very thick, ScoutNavy] (5.5, 1.12) -- (5.5, 0.60) -- (-5.1, 0.60) -- (-5.1, 1.12);
      \node[font=\scriptsize\bfseries\sffamily, ScoutNavy] at (0.2, 0.36) {1. Arco Referencial Externo (\texttt{cod\_depto\_fk} $\longrightarrow$ \underline{\texttt{cod\_depto}})};

      % Arco 2: Auto-referencia MEDICO.mat_superv_fk -> MEDICO.matricula
      \draw[-{Latex[length=2.2mm]}, very thick, ScoutGreen] (7.3, 1.78) -- (7.3, 2.25) -- (1.3, 2.25) -- (1.3, 1.78);
      \node[font=\scriptsize\bfseries\sffamily, ScoutGreen] at (4.3, 2.50) {2. Auto-Referência Recursiva (\texttt{mat\_superv\_fk} $\longrightarrow$ \underline{\texttt{matricula}})};

      % Arco 3: DEPENDENTE.id_prontuario_fk -> PACIENTE.id_prontuario
      \draw[-{Latex[length=2.2mm]}, very thick, ScoutNavy] (2.3, -0.83) -- (2.3, -1.32) -- (-5.2, -1.32) -- (-5.2, -0.83);
      \node[font=\scriptsize\bfseries\sffamily, ScoutNavy] at (-1.45, -1.58) {3. Dependência de Identificação: \texttt{id\_prontuario\_fk} compõe a PK e é FK (Não aceita \texttt{NULL}!)};
    }
  \end{center}
  \vspace{-0.05cm}
  \scriptsize
  $\bullet$ \textbf{1. FK Externa Opcional (\texttt{cod\_depto\_fk}):} Aponta para outra relação; admite \texttt{NULL} quando a lotação é opcional.\\
  $\bullet$ \textbf{2. FK Auto-Referenciada (\texttt{mat\_superv\_fk}):} Aponta para a PK da \textbf{própria relação} (\texttt{MEDICO}), modelando hierarquia.\\
  $\bullet$ \textbf{3. FK em Chave Primária Composta (\texttt{id\_prontuario\_fk}):} Originada de entidade fraca; por integrar a PK, \textbf{jamais aceita \texttt{NULL}}.
\end{frame}

% ------------------------------------------------------------------------------
\begin{frame}{Auditoria Tabular de Tuplas: Verificação Formal de Integridade}
  \vspace{-0.15cm}
  \footnotesize\alert{Avaliação de Operações sobre o Estado da Relação \texttt{MEDICO}:}\par
  \vspace{0.12cm}
  \renewcommand{\arraystretch}{1.42}
  \resizebox{\textwidth}{!}{%
    \begin{tabular}{@{}p{3.5cm}lcp{4.6cm}@{}}
      \toprule
      \textbf{Operação Proposta} & \textbf{Tupla Avaliada $t = \langle v_1, \dots, v_6 \rangle$} & \textbf{Resultado} & \textbf{Justificativa Formal (Codd, 1970)} \\
      \midrule
      1. Inserir médico sem matrícula definida &
      $\langle \text{\textbf{NULL}}, \text{'111'}, \text{'Dr. Rui'}, \text{'CE-10'}, 10, \text{NULL} \rangle$ &
      \textbf{\textcolor{red!80!black}{REJEITADA}} &
      Viola a \textbf{Integridade de Entidade}: nenhum atributo da PK pode ser \texttt{NULL}. \\
      2. Inserir lista de 2 CRMs na mesma célula &
      $\langle 1004, \text{'222'}, \text{'Dr. Ivo'}, \{\text{'CE-1'}, \text{'SP-2'}\}, 10, \text{NULL} \rangle$ &
      \textbf{\textcolor{red!80!black}{REJEITADA}} &
      Viola a \textbf{Integridade de Domínio (1FN)}: exige valores atômicos indivisíveis. \\
      3. Inserir CPF já existente em outra tupla &
      $\langle 1005, \text{'111.222.333-44'}, \text{'Dra. Lia'}, \text{'CE-99'}, 20, \text{NULL} \rangle$ &
      \textbf{\textcolor{red!80!black}{REJEITADA}} &
      Viola a \textbf{Unicidade de Chave}: \texttt{cpf} é Chave Candidata Alternativa ($AK$). \\
      4. Lotar médico no depto. $99$ (inexistente) &
      $\langle 1006, \text{'333'}, \text{'Dr. Gil'}, \text{'CE-55'}, \mathbf{99}, 1001 \rangle$ &
      \textbf{\textcolor{red!80!black}{REJEITADA}} &
      Viola a \textbf{Integridade Referencial}: valor $99 \notin \text{DEPTO}[\underline{\text{cod\_depto}}]$. \\
      5. Inserir médico sem depto. e sem preceptor &
      $\langle 1007, \text{'444'}, \text{'Dra. Eva'}, \text{'CE-77'}, \text{\textbf{NULL}}, \text{\textbf{NULL}} \rangle$ &
      \textbf{\textcolor{ScoutGreen}{ACEITA}} &
      Ambas as Chaves Estrangeiras são opcionais e não integram a PK. \\
      \bottomrule
    \end{tabular}%
  }
\end{frame}

\changecolours{ScoutGreen}
\section[Conclusão e Próximos Passos]{Síntese Conceitual, Exercícios e Ponte para a Aula 10}

% ------------------------------------------------------------------------------
\twocolframe[title={Resumo Executivo da Aula 09 e Próximo Passo},leftcol=7.1cm,rightcol=6.9cm]{%
  \alert{As 4 Regras de Ouro de Codd Consolidadas}\par
  \itemseps[topsepi=0.5mm,itemsepi=1mm,labelsep=2mm,leftmargini=3mm]
  \begin{itemize}
    \item \textbf{1. Domínio Atômico (1FN):} Valores indivisíveis e compatíveis com o tipo do atributo.
    \item \textbf{2. Unicidade de Chave:} Nenhuma duplicata na $PK$ ou nas Chaves Candidatas ($CK$).
    \item \textbf{3. Integridade de Entidade:} $\text{PK} \neq \text{NULL}$ sempre.
    \item \textbf{4. Integridade Referencial:} $\text{FK} \in \text{PK}_{\text{pai}}$ ou $\text{FK} = \text{NULL}$.
  \end{itemize}
}{%
  \alert{Ponte Pedagógica: Aula 10 (27/10/2026)}\par
  \itemseps[topsepi=0.5mm,itemsepi=1mm,labelsep=2mm,leftmargini=3mm]
  \begin{itemize}
    \item \textbf{Agora que dominamos o DER (Aulas 03--08) e o Modelo Relacional (Aula 09):} Como transformar sistematicamente qualquer Diagrama Entidade-Relacionamento em um Esquema Relacional de Tabelas?
    \item \textbf{Aula 10 -- Mapeamento Lógico Relacional (Parte 1):} Algoritmo formal de 7 passos de Elmasri \& Navathe (entidades fortes, entidades fracas e relacionamentos $1:1$ e $1:N$)!
  \end{itemize}
}

\changecolours{ScoutGreen}
\section[Referências]{Referências Bibliográficas Recomendadas}

% ------------------------------------------------------------------------------
\begin{frame}{Referências Bibliográficas}
  \itemseps[topsepi=0.5mm,itemsepi=1mm,labelsep=2mm,leftmargini=3.5mm]
  \begin{itemize}
    \scriptsize
    \item \textbf{CODD, Edgar Frank.} \textit{A Relational Model of Data for Large Shared Data Banks}. Communications of the ACM (CACM), New York, v. 13, n. 6, p. 377--387, June 1970.
    \item \textbf{ELMASRI, Ramez; NAVATHE, Shamkant B.} \textit{Sistemas de Banco de Dados}. 7. ed. São Paulo: Pearson Education do Brasil, 2018. (Capítulo 5: O Modelo de Dados Relacional e as Restrições de Banco de Dados Relacional).
    \item \textbf{HEUSER, Carlos Alberto.} \textit{Projeto de Banco de Dados}. 6. ed. Porto Alegre: Bookman, 2009. (Capítulo 4: Projeto Lógico -- O Modelo Relacional).
    \item \textbf{SILBERSCHATZ, Abraham; KORTH, Henry F.; SUDARSHAN, S.} \textit{Sistema de Banco de Dados}. 7. ed. Rio de Janeiro: Elsevier, 2020. (Capítulo 2: Introdução ao Modelo Relacional).
    \item \textbf{DATE, C. J.} \textit{Introdução a Sistemas de Bancos de Dados}. 8. ed. Rio de Janeiro: Elsevier / Campus, 2004.
  \end{itemize}
  \vspace{0.15cm}
  \begin{center}
    \tiny\color{gray}
    Instituto Federal de Educação, Ciência e Tecnologia do Ceará -- Campus Tauá\\
    Curso Superior de Tecnologia em Telemática
  \end{center}
\end{frame}


\end{document}
